\subsection{Provider-Neutral Persistence Strategy for Product and User Data}
\label{subsec:providerneutrale-persistenzstrategie}

The baseline separates template and development data, runtime configuration,
and product and user data. \texttt{profiles/},
\texttt{project-profile.toml}, \texttt{tools/}, and \texttt{shared/} describe,
generate, or verify the project. \texttt{.env},
\texttt{config/environment.toml}, \texttt{DATABASE\_URL}, and API
configuration instead determine how an instance runs. Documents, planning
states, domain records, media, project files, and histories arise only through
product use and require their own area of responsibility and lifecycle.
Configuration files are not user-data storage; version-controlled template
assets are not runtime assets of the finished product.

A persistence decision distinguishes provider-neutral local file storage,
embedded databases, remote databases, and asset storage. Not every product
requires all four classes. Local files may use JSON, Markdown, or
product-specific formats; an embedded database could, for example, use
SQLite. In contrast, the existing \texttt{database} and \texttt{postgres}
capabilities implement server-side SQL persistence with SQLAlchemy, Alembic,
and PostgreSQL. Asset storage for images, documents, graphics, or binary files
remains separate from structured data. \texttt{local-files},
\texttt{embedded-database}, \texttt{sqlite}, \texttt{object-storage}, and
synchronization are possible extension paths, not implemented capabilities
\parencite{kleiveist2026persistence}.

Each product project must identify the authoritative data source. In a local
application, a local store may be the source of truth; in a server-centered
application, it may be a central database. Offline or local-first systems
additionally require explicit rules for synchronization direction, revisions,
and conflicts; the baseline provides no synchronization engine. The product
decision also covers the schema and versioning strategy, migration, backup and
recovery, and security boundaries.

The template therefore standardizes neither the data itself, a universal
format, nor a general storage provider. It standardizes the boundaries of
responsibility and the decisions that must be documented for handling data.
Storage format and persistence architecture, platform profile and storage
provider, configuration and user data, and asset storage and structured data
storage remain separate decisions.
\beginnerinput{workspace/statement/300_architektur/380}
